\documentclass[11pt]{article}

\usepackage[margin=1in]{geometry}
\usepackage[T1]{fontenc}
\usepackage{lmodern}
\usepackage{amsmath}
\usepackage{booktabs}
\usepackage[version=4]{mhchem}
\usepackage{cleveref}

\Crefname{table}{Table}{Tables}
\crefname{table}{Table}{Tables}
\Crefname{equation}{Equation}{Equations}
\crefname{equation}{Equation}{Equations}

\title{Notation for Aqueous Reaction Chemistry:\\ A Short Worked Survey with \texttt{mhchem}}
\author{A.\ Khoo}
\date{\today}

\begin{document}
\maketitle

\begin{abstract}
This note surveys a handful of elementary results in aqueous and combustion
chemistry, using the \texttt{mhchem} package (version~4) to typeset every
chemical formula and reaction. We move from simple molecular formulae through
polyatomic ions to balanced reaction equations, and close with a small table of
physical properties indexed by chemical species. The aim is expository: the
chemistry is standard, and the emphasis falls on consistent, readable notation
within ordinary academic prose.
\end{abstract}

\section{Molecular formulae in running text}

The most basic use of \ce{} is to set a molecular formula inline. Water,
\ce{H2O}, and carbon dioxide, \ce{CO2}, are the canonical products of complete
hydrocarbon combustion, while methane, \ce{CH4}, is the simplest fuel. Subscript
stoichiometry is handled automatically: glucose is written \ce{C6H12O6}, and
ammonia is \ce{NH3}. When a formula carries net charge, \texttt{mhchem} places
the sign as a superscript without further intervention, so the sulfate ion is
\ce{SO4^2-}, the hydronium ion is \ce{H3O+}, the hydroxide ion is \ce{OH-}, and
the calcium cation is \ce{Ca^2+}. Isotopic and state annotations compose in the
same way; for instance dissolved oxygen may be written \ce{O2(aq)} and solid
calcium carbonate \ce{CaCO3(s)}.

\section{Balanced reactions inline and in display}

A balanced equation may be set inline when it is short. The combustion of
hydrogen, \ce{2 H2 + O2 -> 2 H2O}, and the neutralisation of hydrochloric acid by
sodium hydroxide, \ce{HCl + NaOH -> NaCl + H2O}, are both compact enough to read
comfortably mid-sentence. Reaction arrows, the $+$ separators, and stoichiometric
coefficients are all produced by the \ce{} parser rather than by hand.

Longer or more important reactions are better displayed. The dissolution and
buffering of carbon dioxide in water, central to ocean acidification, proceeds
through carbonic acid and bicarbonate as in \cref{eq:carbonate}:
%
\begin{equation}
  \ce{CO2(aq) + H2O(l) <=> H2CO3(aq) <=> H+ + HCO3^- <=> 2 H+ + CO3^2-}.
  \label{eq:carbonate}
\end{equation}
%
The doubled equilibrium arrows in \cref{eq:carbonate} indicate reversible steps;
each forward step releases a proton, which is why rising dissolved \ce{CO2}
lowers the $\mathrm{pH}$ of seawater. A second illustrative display is the
complete combustion of methane,
%
\begin{equation}
  \ce{CH4(g) + 2 O2(g) -> CO2(g) + 2 H2O(g)},
  \label{eq:methane}
\end{equation}
%
which conserves four hydrogen atoms and, counting the two diatomic oxygen
molecules against the products, four oxygen atoms on each side. \Cref{eq:methane}
is the elementary energy-release reaction behind natural-gas heating.

\section{Thermodynamic quantities}

Reaction enthalpies attach naturally to the displayed equations. For the methane
combustion of \cref{eq:methane} the standard enthalpy of reaction is
$\Delta H^\circ_{\mathrm{r}} \approx -890~\mathrm{kJ\,mol^{-1}}$, the strongly
negative sign marking it as exothermic. In general the standard reaction enthalpy
is obtained from tabulated enthalpies of formation by
\[
  \Delta H^\circ_{\mathrm{r}}
  = \sum_{\text{products}} \nu_i\, \Delta H^\circ_{\mathrm{f},i}
  - \sum_{\text{reactants}} \nu_j\, \Delta H^\circ_{\mathrm{f},j},
\]
where the $\nu$ are stoichiometric coefficients read directly off the balanced
equation. The same bookkeeping applied to \cref{eq:carbonate} explains why the
carbonate equilibria are only mildly temperature-sensitive near ambient
conditions.

\section{A table of species and properties}

\Cref{tab:species} collects a few of the species named above together with their
molar masses and aqueous behaviour. Because the leftmost column is set with
\ce{}, the formulae carry their proper subscripts and charges inside the
\texttt{booktabs} rule structure.

\begin{table}[ht]
  \centering
  \caption{Selected chemical species with approximate molar masses and a
  qualitative note on aqueous behaviour. Molar masses are rounded to one
  decimal place.}
  \label{tab:species}
  \begin{tabular}{l c l}
    \toprule
    Species & Molar mass $(\mathrm{g\,mol^{-1}})$ & Aqueous behaviour \\
    \midrule
    \ce{H2O}      & 18.0  & solvent \\
    \ce{CO2}      & 44.0  & weakly acidic on dissolution \\
    \ce{CH4}      & 16.0  & essentially insoluble \\
    \ce{NaCl}     & 58.4  & fully dissociates \\
    \ce{SO4^2-}   & 96.1  & spectator in most acid--base systems \\
    \ce{HCO3^-}   & 61.0  & amphoteric buffer component \\
    \ce{CaCO3}    & 100.1 & sparingly soluble \\
    \bottomrule
  \end{tabular}
\end{table}

The entries of \cref{tab:species} are consistent with the reactions discussed
earlier: \ce{NaCl}, the product of the neutralisation
\ce{HCl + NaOH -> NaCl + H2O}, dissociates completely, whereas \ce{CaCO3}, formed
when \ce{CO3^2-} from \cref{eq:carbonate} meets \ce{Ca^2+}, is only sparingly
soluble and precipitates as a solid.

\section{Conclusion}

Across \crefrange{eq:carbonate}{eq:methane} and \cref{tab:species}, a single
markup convention, the \ce{} command of \texttt{mhchem}, handles molecular
formulae such as \ce{H2O} and \ce{CO2}, charged species such as \ce{SO4^2-},
balanced reactions such as \ce{2 H2 + O2 -> 2 H2O}, reversible equilibria, and
formulae embedded in a \texttt{booktabs} table. The notation remains uniform
whether a formula appears inline in prose, in a numbered display equation, or in
a tabulated column.

\end{document}
